\bookchapter{Where to Go from Here}{ch:06-where-to-go-from-here}

\begin{chapterintro}

You now know what happens between tapping a button and seeing the result. Here is what to learn next, and what to leave alone for now.

\end{chapterintro}

\section{What to Learn Next}\label{06-where-to-go-from-here:what-to-learn-next}

Each step of the trip has its own book in this series:

\begin{itemize}
\tightlist
\item
  The journey between the two sides: \emph{Networking from Zero}.
\item
  The contract and the back end: \emph{APIs to FastAPI}.
\item
  Where the data lives: \emph{Databases from Zero}.
\item
  The tools every developer uses: \emph{Linux and Git from Zero}.
\item
  Getting a back end onto real servers: \emph{Infrastructure Engineering from Zero}.
\item
  Putting it all together: \emph{The Complete Software Engineering Guide}.
\end{itemize}

For the front end, learn plain HTML, CSS, and JavaScript first, and pick a framework only after that. MDN Web Docs is a free, reliable reference for all three.

\section{What to Skip for Now}\label{06-where-to-go-from-here:what-to-skip-for-now}

\begin{itemize}
\tightlist
\item
  \emph{Choosing the ``best'' framework.} They all do the jobs in Chapters 2 and 3. Learn one well.
\item
  \emph{Scaling to millions of users.} Microservices and Kubernetes solve problems you don't have yet.
\item
  \emph{Other ways to talk.} WebSockets, GraphQL, and server-side rendering are variations on the same trip. Meet them when a project needs them.
\end{itemize}

Whatever you learn next, keep the trip in mind: tap, request, work, response, show. Every new tool is a better way to do one of those steps.
